\documentclass[11pt,aspectratio=169]{beamer}
\usepackage[spanish,es-nodecimaldot]{babel}
\usepackage{fontspec}
\setsansfont{Inter 18pt}[BoldFont={Inter 18pt SemiBold}]
\renewcommand{\familydefault}{\sfdefault}
\usepackage{graphicx}
\usepackage{tikz}
\graphicspath{{sesion2/presentation/images/}{images/}}

\definecolor{tinta}{HTML}{191A20}
\definecolor{azul}{HTML}{224765}
\definecolor{tarjeta}{HTML}{F5F5F5}
\definecolor{imagenA}{HTML}{D9E1F2}
\definecolor{imagenB}{HTML}{DDEBE3}

\setbeamercolor{normal text}{fg=tinta,bg=white}
\setbeamertemplate{navigation symbols}{}
\setbeamertemplate{headline}{}
\setbeamertemplate{footline}{}
\hyphenpenalty=10000
\exhyphenpenalty=10000

% Coordinates follow the 1398x786 canvas used in session 1.
\newcommand{\px}[1]{\dimexpr\paperwidth*#1/1398\relax}
\newcommand{\py}[1]{\dimexpr\paperheight*#1/786\relax}
\newcommand{\mentinode}[6]{%
  \node[anchor=north west,inner sep=0pt,outer sep=0pt,text width=\px{#3},align=left]
    at ([xshift=\px{#1},yshift=-\py{#2}]current page.north west)
    {{\color{tinta}\fontsize{#4}{#5}\selectfont #6}};%
}
\newcommand{\centernode}[6]{%
  \node[anchor=center,inner sep=0pt,outer sep=0pt,text width=\px{#3},align=center]
    at ([xshift=\px{#1},yshift=-\py{#2}]current page.north west)
    {{\color{tinta}\fontsize{#4}{#5}\selectfont #6}};%
}
\newcommand{\panel}[4]{%
  \fill[tarjeta,rounded corners=5pt]
    ([xshift=\px{#1},yshift=-\py{#2}]current page.north west)
    rectangle
    ([xshift=\px{#3},yshift=-\py{#4}]current page.north west);%
}
\newcommand{\flecha}[3]{%
  \draw[->,line width=1pt] ([xshift=\px{#1},yshift=-\py{#3}]current page.north west) -- ([xshift=\px{#2},yshift=-\py{#3}]current page.north west);%
}
% Etiqueta de tiempo de la misión, arriba a la derecha.
\newcommand{\tiempo}[1]{%
  \node[anchor=north east,inner sep=4pt,fill=azul,rounded corners=3pt,text=white]
    at ([xshift=-\px{54},yshift=-\py{60}]current page.north east)
    {\fontsize{8}{10}\selectfont #1};%
}
\newcommand{\fila}[5]{%
  \centernode{330}{#1}{330}{8}{11}{#2}
  \centernode{690}{#1}{150}{8}{11}{#3}
  \centernode{900}{#1}{200}{8}{11}{#4}
  \centernode{1120}{#1}{200}{8}{11}{#5}%
}

% Imagen 8:5 (o la proporción indicada): si el archivo existe se incluye; si no, queda un recuadro de reserva.
% #1 archivo, #2 x, #3 y, #4 ancho, #5 alto de la reserva (lienzo 1398x786).
\newcommand{\imagen}[5]{%
  \IfFileExists{images/#1}{\def\hayimagen{1}}{\IfFileExists{sesion2/presentation/images/#1}{\def\hayimagen{1}}{\def\hayimagen{0}}}%
  \if\hayimagen1
    \node[anchor=north west,inner sep=0] at ([xshift=\px{#2},yshift=-\py{#3}]current page.north west)
      {\includegraphics[width=\px{#4}]{#1}};%
  \else
    \fill[imagenB,rounded corners=5pt] ([xshift=\px{#2},yshift=-\py{#3}]current page.north west)
      rectangle ++(\px{#4},-\py{#5});%
    \node[anchor=center,text=tinta] at ([xshift=\px{#2}+\px{#4}/2,yshift=-\py{#3}-\py{#5}/2]current page.north west)
      {\fontsize{8}{10}\selectfont #1};%
  \fi
}

\definecolor{revisar}{HTML}{F6DCCB}
\newcommand{\pt}[2]{([xshift=\px{#1},yshift=-\py{#2}]current page.north west)}
\newcommand{\caja}[5]{\fill[#5,rounded corners=5pt] \pt{#1}{#2} rectangle \pt{#3}{#4};}

\begin{document}

% 1. Inicio.
\begin{frame}[plain]
  \begin{tikzpicture}[remember picture,overlay]
    \mentinode{54}{247}{650}{30}{36}{IA en el día a día}
    \mentinode{55}{345}{620}{11}{16}{Texto, correos, documentos y herramientas que ya conocemos.}
    \imagen{portada-herramientas.png}{700}{300}{640}{400}
  \end{tikzpicture}
\end{frame}

% 2. Juego de arranque.
\begin{frame}[plain]
  \begin{tikzpicture}[remember picture,overlay]
    \mentinode{54}{170}{1280}{19.5}{24}{¿Qué palabra sigue?}
    \mentinode{55}{265}{1280}{9}{14}{Escribe en la ficha la primera palabra que se te ocurra. Tienes 10 segundos.}
    \panel{54}{340}{453}{580}
    \panel{499}{340}{899}{580}
    \panel{945}{340}{1344}{580}
    \centernode{254}{390}{340}{9}{13}{Ronda 1}
    \centernode{699}{390}{340}{9}{13}{Ronda 2}
    \centernode{1145}{390}{340}{9}{13}{Ronda 3}
    \centernode{254}{480}{340}{11}{15}{«En la selva amazónica viven muchos…»}
    \centernode{699}{480}{340}{11}{15}{«Mi comida favorita es el…»}
    \centernode{1145}{480}{340}{11}{15}{«La primera panadería de San Juan Bosco la abrió…»}
    \mentinode{55}{640}{1280}{8}{14}{Mano arriba según tu respuesta. ¿Coincidimos? ¿Alguien sabía la tercera?}
  \end{tikzpicture}
\end{frame}

% 3. Cómo funciona: ciclo de uso (guía del instructor).
\begin{frame}[plain]
  \begin{tikzpicture}[remember picture,overlay]
    \mentinode{54}{130}{1280}{19.5}{24}{Cómo funciona}
    \mentinode{55}{215}{1280}{9}{14}{Escribe eligiendo una palabra probable tras otra. \textcolor{azul}{Probable no es lo mismo que verdadero.}}
    \caja{90}{290}{440}{390}{imagenA}
    \caja{525}{290}{875}{390}{imagenA}
    \caja{960}{290}{1310}{390}{imagenA}
    \caja{500}{445}{900}{545}{revisar}
    \centernode{265}{340}{320}{10}{14}{1. Escribo una instrucción}
    \centernode{700}{340}{320}{10}{14}{2. La IA identifica patrones y contexto}
    \centernode{1135}{340}{320}{10}{14}{3. Genera una respuesta}
    \centernode{700}{495}{370}{10}{14}{4. La persona revisa, corrige y decide si usarla}
    \draw[->,line width=1pt] \pt{440}{340} -- \pt{525}{340};
    \draw[->,line width=1pt] \pt{875}{340} -- \pt{960}{340};
    \draw[->,line width=1pt] \pt{1135}{390} |- \pt{900}{495};
    \draw[->,line width=1pt] \pt{500}{495} -| \pt{265}{390};
    \caja{54}{600}{1344}{690}{azul}
    \centernode{699}{645}{1240}{9}{14}{\color{white}\textbf{Atención:} la respuesta puede sonar muy segura aunque tenga errores. La revisión humana es obligatoria.}
  \end{tikzpicture}
\end{frame}

% 4. Lo que hace y lo que no hace.
\begin{frame}[plain]
  \begin{tikzpicture}[remember picture,overlay]
    \mentinode{54}{170}{1280}{19.5}{24}{Lo que hace y lo que no hace}
    \panel{54}{290}{661}{620}
    \panel{737}{290}{1344}{620}
    \mentinode{95}{335}{540}{12}{16}{Sí hace}
    \mentinode{95}{405}{540}{9}{13}{Escribe, resume y traduce.}
    \mentinode{95}{460}{540}{9}{13}{Usa lo que dijiste en el mismo chat.}
    \mentinode{95}{515}{540}{9}{13}{Mejora si le pides cambios.}
    \mentinode{778}{335}{540}{12}{16}{No hace}
    \mentinode{778}{405}{540}{9}{13}{No comprueba si es verdad.}
    \mentinode{778}{460}{540}{9}{13}{No sabe lo más reciente.}
    \mentinode{778}{515}{540}{9}{13}{No recuerda tus otros chats.}
    \mentinode{778}{570}{540}{9}{13}{No responde siempre igual.}
    \mentinode{55}{670}{1280}{8}{14}{Prueba: pide lo mismo dos veces. ¿Salió igual?}
  \end{tikzpicture}
\end{frame}

% 5. Demostración en vivo.
\begin{frame}[plain]
  \begin{tikzpicture}[remember picture,overlay]
    \mentinode{54}{220}{1280}{19.5}{24}{Mejoremos una petición entre todos}
    \panel{54}{340}{1344}{430}
    \centernode{699}{385}{1190}{10}{14}{«Háblame de la Amazonía»}
    \panel{54}{460}{453}{620}
    \panel{499}{460}{899}{620}
    \panel{945}{460}{1344}{620}
    \centernode{254}{505}{340}{9}{13}{Ronda 1}
    \centernode{699}{505}{340}{9}{13}{Ronda 2}
    \centernode{1145}{505}{340}{9}{13}{Ronda 3}
    \centernode{254}{565}{340}{11}{15}{¿Para quién es?}
    \centernode{699}{565}{340}{11}{15}{¿Qué necesitamos?}
    \centernode{1145}{565}{340}{11}{15}{¿Cómo la queremos?}
    \mentinode{55}{660}{1280}{8}{11}{Dicten en voz alta qué agregar. Probamos después de cada ronda.}
  \end{tikzpicture}
\end{frame}

% 6. Fórmula de un buen prompt (guía del instructor).
\begin{frame}[plain]
  \begin{tikzpicture}[remember picture,overlay]
    \mentinode{54}{110}{1280}{19.5}{24}{Fórmula de un buen prompt}
    \centernode{699}{230}{1290}{13}{17}{Rol + Tarea + Contexto + Formato + Tono + Límite}
    \panel{54}{275}{1344}{539}
    \foreach \y in {319,363,407,451,495}
      \draw[white,line width=1.5pt] \pt{54}{\y} -- \pt{1344}{\y};
    \mentinode{110}{286}{250}{9}{13}{\textbf{Rol}}
    \mentinode{400}{286}{900}{9}{13}{¿Desde qué perspectiva debe responder?}
    \mentinode{110}{330}{250}{9}{13}{\textbf{Tarea}}
    \mentinode{400}{330}{900}{9}{13}{¿Qué debe hacer exactamente?}
    \mentinode{110}{374}{250}{9}{13}{\textbf{Contexto}}
    \mentinode{400}{374}{900}{9}{13}{¿Qué información necesita saber?}
    \mentinode{110}{418}{250}{9}{13}{\textbf{Formato}}
    \mentinode{400}{418}{900}{9}{13}{¿Cómo quiero recibir la respuesta?}
    \mentinode{110}{462}{250}{9}{13}{\textbf{Tono}}
    \mentinode{400}{462}{900}{9}{13}{¿Formal, amable, directo, institucional?}
    \mentinode{110}{506}{250}{9}{13}{\textbf{Límite}}
    \mentinode{400}{506}{900}{9}{13}{¿Cuántas palabras, puntos o párrafos?}
    \caja{54}{560}{1344}{690}{imagenB}
    \centernode{699}{625}{1220}{8}{13}{«Actúa como guía turístico. Necesito un texto que invite a conocer San Juan Bosco, en Morona Santiago. Es para familias que visitan el cantón por primera vez. Escríbelo como una lista de cinco puntos, con tono amable, en máximo 80 palabras y sin inventar lugares ni horarios»}
    \mentinode{55}{715}{1280}{8}{11}{¿Dónde está cada elemento en este ejemplo?}
  \end{tikzpicture}
\end{frame}

% 7. Plantilla para copiar (guía del instructor).
\begin{frame}[plain]
  \begin{tikzpicture}[remember picture,overlay]
    \tiempo{5 min}
    \mentinode{54}{120}{1280}{19.5}{24}{Plantilla para copiar}
    \panel{54}{195}{1344}{480}
    \mentinode{100}{220}{1200}{9}{15.5}{Actúa como \textcolor{azul}{[rol]}.\\Necesito que \textcolor{azul}{[tarea]}.\\El contexto es \textcolor{azul}{[información clave]}.\\El resultado debe estar en formato \textcolor{azul}{[correo / tabla / lista / informe]}.\\Usa un tono \textcolor{azul}{[formal / amable / institucional / directo]}.\\Limita la respuesta a \textcolor{azul}{[número de palabras / párrafos / puntos]}.}
    \caja{54}{505}{1344}{680}{imagenB}
    \mentinode{100}{525}{1200}{9}{13}{\textbf{Pruébala ahora con las notas de la ficha}}
    \mentinode{100}{575}{1200}{8}{13}{«Actúa como coordinador de un club escolar. Necesito que conviertas estas notas en una lista de tareas. El contexto es una reunión para organizar nuestro puesto en la Feria Comunitaria. El resultado debe estar en una tabla. Usa un tono claro y amable. Limita cada tarea a una frase.»}
    \mentinode{55}{715}{1280}{8}{11}{La plantilla está en la ficha: úsala en todos los niveles.}
  \end{tikzpicture}
\end{frame}

% 7. Taller: mapa de misiones.
\begin{frame}[plain]
  \begin{tikzpicture}[remember picture,overlay]
    \mentinode{54}{200}{1280}{19.5}{24}{Taller: cinco misiones}
    \panel{54}{320}{296}{540}
    \panel{316}{320}{558}{540}
    \panel{578}{320}{820}{540}
    \panel{840}{320}{1082}{540}
    \panel{1102}{320}{1344}{540}
    \centernode{175}{370}{210}{9}{13}{Nivel 1}
    \centernode{437}{370}{210}{9}{13}{Nivel 2}
    \centernode{699}{370}{210}{9}{13}{Nivel 3}
    \centernode{961}{370}{210}{9}{13}{Nivel 4}
    \centernode{1223}{370}{210}{9}{13}{Final}
    \centernode{175}{440}{210}{12}{15}{Compara}
    \centernode{437}{440}{210}{12}{15}{Dale un rol}
    \centernode{699}{440}{210}{12}{15}{Formato, tono y límite}
    \centernode{961}{440}{210}{12}{15}{Conversa}
    \centernode{1223}{440}{210}{12}{15}{Batalla}
    \centernode{175}{505}{210}{8}{11}{8 min}
    \centernode{437}{505}{210}{8}{11}{8 min}
    \centernode{699}{505}{210}{8}{11}{8 min}
    \centernode{961}{505}{210}{8}{11}{8 min}
    \centernode{1223}{505}{210}{8}{11}{12 min}
    \mentinode{55}{590}{1280}{8}{14}{Cada uno en su computador. Copia tu prompt en la ficha y anota lo que observaste.}
    \mentinode{55}{635}{1280}{8}{14}{¿Terminaste antes? Haz la opción extra de cada nivel.}
  \end{tikzpicture}
\end{frame}

% 8. Nivel 1.
\begin{frame}[plain]
  \begin{tikzpicture}[remember picture,overlay]
    \tiempo{8 min}
    \mentinode{54}{200}{1280}{19.5}{24}{Nivel 1: misma pregunta, distintas IA}
    \panel{54}{310}{1344}{430}
    \centernode{699}{370}{1190}{9}{14}{«Explica qué es la biodiversidad a un estudiante de 15 años en tres frases y da un ejemplo de la Amazonía»}
    \mentinode{55}{480}{1280}{9}{14}{Envíala a dos asistentes: ChatGPT, Gemini o Claude.}
    \mentinode{55}{540}{1280}{8}{14}{\textbullet\ ¿Cuál respuesta es más clara?}
    \mentinode{55}{585}{1280}{8}{14}{\textbullet\ ¿En qué se diferencian?}
    \mentinode{55}{630}{1280}{8}{14}{\textbullet\ ¿Hay algún dato que deberías comprobar?}
  \end{tikzpicture}
\end{frame}

% 9. Nivel 2.
\begin{frame}[plain]
  \begin{tikzpicture}[remember picture,overlay]
    \tiempo{8 min}
    \mentinode{54}{200}{1280}{19.5}{24}{Nivel 2: dale un rol}
    \mentinode{55}{295}{1280}{9}{14}{Escribe tu propio prompt para dos opciones. Empieza con «Actúa como…».}
    \panel{54}{370}{453}{620}
    \panel{499}{370}{899}{620}
    \panel{945}{370}{1344}{620}
    \centernode{254}{420}{340}{12}{16}{A. Guía turístico}
    \centernode{699}{420}{340}{12}{16}{B. Profesor}
    \centernode{1145}{420}{340}{12}{16}{C. Contador}
    \centernode{254}{525}{340}{8}{13}{Un día en San Juan Bosco para jóvenes visitantes, sin inventar lugares.}
    \centernode{699}{525}{340}{8}{13}{La erosión del suelo con un ejemplo de agricultura.}
    \centernode{1145}{525}{340}{8}{13}{Cómo calcula su ganancia una familia que vende queso.}
    \mentinode{55}{660}{1280}{8}{11}{¿Qué cambió en la respuesta al darle un rol?}
  \end{tikzpicture}
\end{frame}

% 10. Nivel 3.
\begin{frame}[plain]
  \begin{tikzpicture}[remember picture,overlay]
    \tiempo{8 min}
    \mentinode{54}{200}{1280}{19.5}{24}{Nivel 3: formato, tono y límite}
    \mentinode{55}{295}{1280}{9}{14}{Dos opciones. Pide un formato, un tono y un límite (palabras, puntos o «no inventes»).}
    \panel{54}{370}{453}{620}
    \panel{499}{370}{899}{620}
    \panel{945}{370}{1344}{620}
    \centernode{254}{420}{340}{12}{16}{A. Tabla}
    \centernode{699}{420}{340}{12}{16}{B. Plan de negocio}
    \centernode{1145}{420}{340}{12}{16}{C. Resumen}
    \centernode{254}{525}{340}{8}{13}{Actividades de un negocio familiar, dificultades y mejoras.}
    \centernode{699}{525}{340}{8}{13}{Jugos en la Feria: producto, clientes, costos, precio y riesgo.}
    \centernode{1145}{525}{340}{8}{13}{El texto de la ficha en cuatro ideas, sin añadir nada.}
    \mentinode{55}{660}{1280}{8}{11}{¿Respetó formato, tono y límite? ¿Agregó algo que no pediste?}
  \end{tikzpicture}
\end{frame}

% 11. Nivel 4: conversar e iterar.
\begin{frame}[plain]
  \begin{tikzpicture}[remember picture,overlay]
    \tiempo{8 min}
    \mentinode{54}{200}{1280}{19.5}{24}{Nivel 4: conversa con la IA}
    \mentinode{55}{295}{1280}{9}{14}{La primera respuesta es un borrador. Toma una del nivel 2 o 3 y envía tres mensajes más.}
    \panel{54}{370}{453}{580}
    \panel{499}{370}{899}{580}
    \panel{945}{370}{1344}{580}
    \centernode{254}{420}{340}{9}{13}{Mensaje 1}
    \centernode{699}{420}{340}{9}{13}{Mensaje 2}
    \centernode{1145}{420}{340}{9}{13}{Mensaje 3}
    \centernode{254}{500}{340}{10}{14}{«Hazlo más corto y sencillo»}
    \centernode{699}{500}{340}{10}{14}{«Agrega un ejemplo de San Juan Bosco»}
    \centernode{1145}{500}{340}{10}{14}{«¿Qué datos de tu respuesta debo comprobar?»}
    \mentinode{55}{630}{1280}{8}{14}{¿Mejoró? ¿El ejemplo es real?}
    \mentinode{55}{675}{1280}{8}{11}{Extra: dale un modelo. «Escribe otro texto con el mismo estilo que este: …»}
  \end{tikzpicture}
\end{frame}

% 12. Batalla de prompts.
\begin{frame}[plain]
  \begin{tikzpicture}[remember picture,overlay]
    \tiempo{6 min + votación}
    \mentinode{54}{200}{1280}{19.5}{24}{Batalla de prompts}
    \panel{54}{300}{1344}{420}
    \centernode{699}{360}{1190}{9}{14}{Consigue un mensaje de máximo 50 palabras que invite a jóvenes de otra ciudad a visitar San Juan Bosco.}
    \panel{54}{460}{354}{580}
    \panel{384}{460}{684}{580}
    \panel{714}{460}{1014}{580}
    \panel{1044}{460}{1344}{580}
    \centernode{204}{520}{260}{11}{15}{Claro}
    \centernode{534}{520}{260}{11}{15}{Atractivo}
    \centernode{864}{520}{260}{11}{15}{50 palabras o menos}
    \centernode{1194}{520}{260}{11}{15}{Sin inventar}
    \mentinode{55}{630}{1280}{8}{14}{Mostramos cuatro resultados y votamos a mano alzada. ¿Qué tenía el prompt ganador?}
  \end{tikzpicture}
\end{frame}

% 13. Cazadores de alucinaciones.
\begin{frame}[plain]
  \begin{tikzpicture}[remember picture,overlay]
    \tiempo{12 min}
    \mentinode{54}{170}{560}{19.5}{24}{Cazadores de alucinaciones}
    \mentinode{55}{345}{560}{8}{13}{Datos que la IA inventa o equivoca, aunque suene segura.}
    \mentinode{55}{445}{560}{8}{13}{1. Pide tres datos de San Juan Bosco.}
    \mentinode{55}{490}{560}{8}{13}{2. Búscalos en una fuente oficial.}
    \mentinode{55}{535}{560}{8}{13}{3. Compara y anota en la ficha.}
    \mentinode{55}{610}{560}{7}{11}{Ej.: año de cantonización, habitantes, ríos. Fuentes: municipio, INEC.}
    \imagen{cazar-alucinacion.png}{650}{200}{690}{431}
  \end{tikzpicture}
\end{frame}

% 14. Verificación y uso responsable.
\begin{frame}[plain]
  \begin{tikzpicture}[remember picture,overlay]
    \mentinode{54}{170}{560}{19.5}{24}{Antes de usar una respuesta}
    \mentinode{55}{335}{560}{10}{14}{Comprueba}
    \mentinode{55}{380}{560}{8}{13}{¿Responde lo que pediste?}
    \mentinode{55}{420}{560}{8}{13}{¿Hay datos o fechas por verificar?}
    \mentinode{55}{490}{560}{10}{14}{Cuida}
    \mentinode{55}{535}{560}{8}{13}{No compartas datos privados.}
    \mentinode{55}{575}{560}{8}{13}{No copies como propio un texto de IA.}
    \mentinode{55}{615}{560}{8}{13}{Indica cuándo usaste IA.}
    \imagen{comprobar-cuidar.png}{650}{200}{690}{431}
  \end{tikzpicture}
\end{frame}

% 15. Estaciones de ofimática.
\begin{frame}[plain]
  \begin{tikzpicture}[remember picture,overlay]
    \mentinode{54}{200}{1280}{19.5}{24}{Estaciones de ofimática}
    \panel{54}{310}{453}{500}
    \panel{499}{310}{899}{500}
    \panel{945}{310}{1344}{500}
    \centernode{254}{375}{340}{13}{17}{Correo}
    \centernode{699}{375}{340}{13}{17}{Documento}
    \centernode{1145}{375}{340}{13}{17}{Hoja de cálculo}
    \centernode{254}{450}{340}{8}{11}{12 min}
    \centernode{699}{450}{340}{8}{11}{10 min}
    \centernode{1145}{450}{340}{8}{11}{13 min}
    \mentinode{55}{555}{1280}{8}{14}{¿Tienes Gemini en Gmail, Docs o Sheets? Úsalo.}
    \mentinode{55}{600}{1280}{8}{14}{Si no: pide el texto o la fórmula a un asistente de IA y cópialo en el programa.}
    \mentinode{55}{660}{1280}{8}{11}{Microsoft 365 Copilot hace lo mismo en Word, Excel y Outlook.}
  \end{tikzpicture}
\end{frame}

% 16. Correo y documento.
\begin{frame}[plain]
  \begin{tikzpicture}[remember picture,overlay]
    \mentinode{54}{200}{1280}{19.5}{24}{Correo y documento}
    \panel{54}{310}{661}{640}
    \panel{737}{310}{1344}{640}
    \centernode{358}{360}{545}{12}{16}{Correo}
    \centernode{358}{430}{545}{8}{14}{Pide información a una institución con un correo formal.}
    \centernode{358}{500}{545}{8}{14}{Guárdalo como borrador: no lo envíes.}
    \centernode{358}{570}{545}{8}{14}{Resume en tres puntos el hilo de la ficha.}
    \centernode{1041}{360}{545}{12}{16}{Documento}
    \centernode{1041}{430}{545}{8}{14}{Escribe una invitación a la Feria Comunitaria.}
    \centernode{1041}{500}{545}{8}{14}{Pide la traducción al inglés.}
    \centernode{1041}{570}{545}{8}{14}{Revisa que conserve el significado.}
    \mentinode{55}{675}{1280}{8}{11}{Borrador \(\rightarrow\) revisar \(\rightarrow\) editar}
  \end{tikzpicture}
\end{frame}

% 17. Hoja de cálculo.
\begin{frame}[plain]
  \begin{tikzpicture}[remember picture,overlay]
    \tiempo{13 min}
    \mentinode{54}{150}{1280}{19.5}{24}{Hoja de cálculo: puesto de jugos}
    \mentinode{55}{245}{1280}{8}{14}{Copia la tabla, pide a la IA las fórmulas, comprueba el resultado y crea un gráfico de gastos.}
    \panel{160}{300}{1238}{624}
    \foreach \y in {344,384,424,464,504,544,584}
      \draw[white,line width=1pt] ([xshift=\px{160},yshift=-\py{\y}]current page.north west) -- ([xshift=\px{1238},yshift=-\py{\y}]current page.north west);
    \fila{322}{\textbf{Producto}}{\textbf{Cantidad}}{\textbf{Precio}}{\textbf{Subtotal}}
    \fila{364}{Naranjas}{30}{\$0,15}{?}
    \fila{404}{Azúcar (kg)}{2}{\$1,10}{?}
    \fila{444}{Vasos}{50}{\$0,04}{?}
    \fila{484}{Hielo (fundas)}{2}{\$1,00}{?}
    \fila{524}{\textbf{Total de gastos}}{}{}{?}
    \fila{564}{Ingresos (jugos vendidos)}{40}{\$0,50}{?}
    \fila{604}{\textbf{Ganancia}}{}{}{?}
    \mentinode{55}{660}{1280}{8}{11}{Cifras ficticias. La IA sugiere la fórmula; tú compruebas el cálculo.}
  \end{tikzpicture}
\end{frame}

% 18. Avance del proyecto final.
\begin{frame}[plain]
  \begin{tikzpicture}[remember picture,overlay]
    \mentinode{54}{120}{560}{19.5}{24}{Avance del proyecto final}
    \mentinode{55}{290}{560}{8}{13}{En equipo; cada uno con un rol.}
    \mentinode{55}{355}{560}{9}{14}{1. Ordenar \textcolor{azul}{\fontsize{7}{9}\selectfont 12 min}}
    \mentinode{55}{393}{560}{7}{11}{Problema, a quién afecta, qué sabemos y qué falta.}
    \mentinode{55}{440}{560}{9}{14}{2. Poner a prueba \textcolor{azul}{\fontsize{7}{9}\selectfont 8 min}}
    \mentinode{55}{478}{560}{7}{11}{Preguntas de jurado y de entrevista.}
    \mentinode{55}{525}{560}{9}{14}{3. Presentar \textcolor{azul}{\fontsize{7}{9}\selectfont 12 min}}
    \mentinode{55}{563}{560}{7}{11}{Cuatro diapositivas, una por integrante.}
    \mentinode{55}{610}{560}{9}{14}{4. Intercambio \textcolor{azul}{\fontsize{7}{9}\selectfont 8 min}}
    \mentinode{55}{648}{560}{7}{11}{Con otro equipo: una estrella y un deseo.}
    \imagen{equipo-proyecto.png}{650}{200}{690}{431}
  \end{tikzpicture}
\end{frame}

% 19. Presentación del proyecto.
\begin{frame}[plain]
  \begin{tikzpicture}[remember picture,overlay]
    \tiempo{12 min}
    \mentinode{54}{200}{1280}{19.5}{24}{Su proyecto en cuatro diapositivas}
    \panel{54}{320}{354}{470}
    \panel{384}{320}{684}{470}
    \panel{714}{320}{1014}{470}
    \panel{1044}{320}{1344}{470}
    \centernode{204}{395}{260}{10}{14}{1. Problema}
    \centernode{534}{395}{260}{10}{14}{2. A quién afecta}
    \centernode{864}{395}{260}{10}{14}{3. Lo que averiguamos}
    \centernode{1194}{395}{260}{10}{14}{4. Próximo paso}
    \mentinode{55}{530}{1280}{8}{14}{Cada integrante hace una con ayuda de la IA o de Gemini en Google Slides.}
    \mentinode{55}{575}{1280}{8}{14}{El coordinador las une y revisa que no haya datos inventados.}
    \mentinode{55}{640}{1280}{8}{11}{Pocas palabras por diapositiva: el resto lo dicen ustedes.}
  \end{tikzpicture}
\end{frame}

% 20. Puesta en común de proyectos.
\begin{frame}[plain]
  \begin{tikzpicture}[remember picture,overlay]
    \mentinode{54}{280}{1280}{19.5}{24}{Cada equipo: 90 segundos}
    \panel{54}{400}{354}{550}
    \panel{384}{400}{684}{550}
    \panel{714}{400}{1014}{550}
    \panel{1044}{400}{1344}{550}
    \centernode{204}{475}{260}{11}{15}{Problema}
    \centernode{534}{475}{260}{11}{15}{A quién afecta}
    \centernode{864}{475}{260}{11}{15}{Lo que averiguamos}
    \centernode{1194}{475}{260}{11}{15}{Próximo paso}
  \end{tikzpicture}
\end{frame}

% 21. Cierre y tarea.
\begin{frame}[plain]
  \begin{tikzpicture}[remember picture,overlay]
    \mentinode{54}{170}{1280}{19.5}{24}{Antes de salir}
    \mentinode{55}{270}{1280}{9}{14}{Boleto de salida en la ficha: una ayuda, una precaución y una duda.}
    \mentinode{54}{360}{1280}{14}{18}{Para la próxima semana}
    \panel{54}{430}{1344}{535}
    \centernode{240}{483}{310}{9}{14}{Tarea escolar con IA}
    \flecha{410}{455}{483}
    \centernode{676}{483}{330}{9}{14}{Qué pediste y qué cambiaste}
    \flecha{880}{925}{483}
    \centernode{1128}{483}{340}{9}{14}{Qué comprobaste}
    \mentinode{55}{575}{1280}{8}{14}{Proyecto: consigue la información que quedó en «qué nos falta averiguar».}
    \mentinode{55}{640}{1280}{8}{11}{Guarda tu ficha y entrégala por el medio que indique el instructor.}
  \end{tikzpicture}
\end{frame}

% 22. Reserva: prompt inverso.
\begin{frame}[plain]
  \begin{tikzpicture}[remember picture,overlay]
    \tiempo{10 min}
    \mentinode{55}{110}{1280}{8}{11}{\textcolor{azul}{Actividad de reserva}}
    \mentinode{54}{200}{1280}{19.5}{24}{Prompt inverso: ¿qué pidieron?}
    \panel{54}{290}{1344}{560}
    \mentinode{110}{320}{1180}{9}{16}{1. Lleva botas de caucho y ropa de manga larga.\\2. Usa repelente contra mosquitos.\\3. Lleva agua suficiente.\\4. No te salgas del sendero.\\5. Avisa a alguien a dónde vas.}
    \mentinode{55}{610}{1280}{8}{14}{Escribe en la ficha el prompt que crees que produjo esta respuesta y pruébalo.}
    \mentinode{55}{655}{1280}{8}{14}{Gana el resultado más parecido. ¿Qué claves tenía?}
  \end{tikzpicture}
\end{frame}

% 23. Reserva: que la IA te entreviste.
\begin{frame}[plain]
  \begin{tikzpicture}[remember picture,overlay]
    \tiempo{10 min}
    \mentinode{55}{110}{1280}{8}{11}{\textcolor{azul}{Actividad de reserva}}
    \mentinode{54}{200}{1280}{19.5}{24}{Que la IA les haga preguntas}
    \panel{54}{300}{1344}{450}
    \centernode{699}{375}{1190}{9}{14}{«Queremos mejorar nuestro proyecto: [idea]. Antes de darnos sugerencias, haznos cinco preguntas, una por una, y espera nuestra respuesta. Al final, danos tres sugerencias.»}
    \mentinode{55}{510}{1280}{9}{14}{En equipo: respondan entre todos.}
    \mentinode{55}{570}{1280}{8}{14}{¿Qué pregunta no supieron responder? Agréguenla a «qué nos falta averiguar».}
    \mentinode{55}{640}{1280}{8}{11}{Truco: pedir que la IA pregunte primero da respuestas más útiles.}
  \end{tikzpicture}
\end{frame}

\end{document}
